\documentclass{article}
\usepackage{amsmath,amssymb,amsthm}
\usepackage{algorithm,algorithmic}
\usepackage{hyperref}
\usepackage{bm}
\newtheorem{theorem}{Theorem}
\newtheorem{lemma}{Lemma}
\newtheorem{assumption}{Assumption}
\begin{document}

\section*{Algorithm Name}
\textbf{AdaGrad for scalar (or per‑coordinate) updates with a running sum of gradients.}

\section*{Mathematical Formulation}
Let $f:\mathbb{R}^d\rightarrow\mathbb{R}$ be the objective function.
Denote by $g_k\in\mathbb{R}^d$ a (possibly stochastic) gradient evaluated at the iterate $w_k$:
\[
g_k = \nabla f(w_k) \quad\text{(deterministic case)}\qquad\text{or}\qquad 
\mathbb{E}[g_k\mid w_k]=\nabla f(w_k).
\]

Define the cumulative sum of squared gradients (element‑wise):
\[
G_k \;:=\; \sum_{i=1}^{k} g_i\odot g_i\in\mathbb{R}^d,\qquad 
\sqrt{G_k} \;:=\; \big(\sqrt{[G_k]_j}\big)_{j=1}^d .
\]

The learning‑rate vector at iteration $k$ is
\[
\eta_k \;:=\; \frac{\alpha}{\sqrt{G_k}+\varepsilon\mathbf 1},
\]
where $\alpha>0$ is a global stepsize hyperparameter and $\varepsilon>0$ is a small constant added for numerical stability.

The update rule is
\[
\boxed{\,w_{k+1}= w_k - \eta_k\odot g_k\,}.
\]

For the scalar (one‑dimensional) case the notation collapses to
\[
A_k = \sum_{i=1}^{k} g_i,\qquad 
\eta_k = \frac{\alpha}{\sqrt{\sum_{i=1}^{k} g_i^{2}}+\varepsilon},
\qquad 
w_{k+1}= w_k - \eta_k g_k .
\]

\section*{Pseudocode}
\begin{algorithm}[H]
\caption{Scalar AdaGrad (running sum of gradients)}\label{alg:adagrad}
\begin{algorithmic}[1]
\REQUIRE Initial point $w_0\in\mathbb{R}$, global stepsize $\alpha>0$, stability term $\varepsilon>0$,
        maximum number of iterations $T$.
\STATE $A_0 \leftarrow 0$ \COMMENT{cumulative sum of gradients}
\FOR{$k=0,\dots,T-1$}
    \STATE Compute (stochastic) gradient $g_k = \nabla f(w_k)$
    \STATE $A_{k+1} \leftarrow A_k + g_k$ \COMMENT{running sum}
    \STATE $\eta_k \leftarrow \displaystyle\frac{\alpha}{\sqrt{A_{k+1}^{2}}+\varepsilon}$
    \STATE $w_{k+1} \leftarrow w_k - \eta_k\, g_k$
\ENDFOR
\RETURN $w_T$
\end{algorithmic}
\end{algorithm}

\section*{Dry Run Example}
Consider the simple quadratic $f(w)=(w-3)^2$.
Hence $\nabla f(w)=2(w-3)$.  
Choose $\alpha=0.1$, $\varepsilon=10^{-8}$ and start from $w_0=0$.

\begin{center}
\begin{tabular}{c|c|c|c|c}
$k$ & $w_k$ & $g_k=2(w_k-3)$ & $A_{k+1}$ & $w_{k+1}=w_k-\eta_k g_k$\\ \hline
0 & $0$ & $-6$ & $-6$ & $\displaystyle 0-\frac{0.1}{\sqrt{36}}\cdot(-6)=0+0.1=0.1$\\[4pt]
1 & $0.1$ & $-5.8$ & $-11.8$ & $\displaystyle 0.1-\frac{0.1}{\sqrt{11.8^{2}}}\cdot(-5.8)
               =0.1+0.1\frac{5.8}{11.8}\approx0.1492$\\[4pt]
2 & $0.1492$ & $-5.7016$ & $-17.5016$ &
 $\displaystyle 0.1492-\frac{0.1}{\sqrt{17.5016^{2}}}\cdot(-5.7016)
 \approx0.1492+0.1\frac{5.7016}{17.5016}\approx0.1825$
\end{tabular}
\end{center}

Objective values:
\[
f(w_0)=9,\quad f(w_1)=(0.1-3)^2\approx8.41,\quad
f(w_2)\approx(0.1492-3)^2\approx8.13,\quad
f(w_3)\approx(0.1825-3)^2\approx7.94 .
\]

The example illustrates that the step size shrinks as the accumulated squared gradients grow.

\newpage
\section*{Assumptions}
\begin{assumption}[Smoothness]\label{ass:smooth}
$f:\mathbb{R}^d\rightarrow\mathbb{R}$ is $L$‑smooth, i.e.,
\[
\|\nabla f(x)-\nabla f(y)\|\le L\|x-y\|\qquad\forall x,y\in\mathbb{R}^d .
\]
Equivalently,
\[
f(y)\le f(x)+\langle\nabla f(x),y-x\rangle+\frac{L}{2}\|y-x\|^{2}.
\]
\end{assumption}
\begin{assumption}[Bounded (second) moments of stochastic gradients]\label{ass:bounded}
For every iteration $k$,
\[
\mathbb{E}[g_k\mid w_k]=\nabla f(w_k),\qquad 
\mathbb{E}\big[\|g_k\|^{2}\mid w_k\big]\le G_{\max }^{2},
\]
for some constant $G_{\max}>0$.
\end{assumption}
\begin{assumption}[Lower bounded objective]\label{ass:lower}
There exists $f^{\star}>-\infty$ such that $f(w)\ge f^{\star}$ for all $w$.
\end{assumption}

All constants $\alpha>0$, $\varepsilon>0$, $L>0$, $G_{\max}>0$ are fixed throughout the analysis.

\section*{Main Convergence Theorem}
\begin{theorem}[Non‑convex convergence of scalar AdaGrad]\label{thm:main}
Let $\{w_k\}_{k\ge0}$ be generated by Algorithm~\ref{alg:adagrad} under
Assumptions~\ref{ass:smooth}–\ref{ass:lower}.  
Define $T\in\mathbb{N}$ and let $G_T:=\sum_{k=0}^{T-1}\mathbb{E}\|g_k\|^{2}$.
Then
\[
\frac{1}{T}\sum_{k=0}^{T-1}\mathbb{E}\big[\|\nabla f(w_k)\|^{2}\big]
\;\le\;
\underbrace{\frac{2\big(f(w_{0})-f^{\star}\big)}{\alpha\sqrt{T}}}_{\text{optimization term}}
\;+\;
\underbrace{\frac{L\alpha G_{\max}}{\sqrt{T}}}_{\text{variance term}} .
\]
Consequently,
\[
\min_{0\le k<T}\mathbb{E}\big[\|\nabla f(w_k)\|^{2}\big]
= O\!\left(\frac{1}{\sqrt{T}}\right).
\]
\end{theorem}

\section*{Theoretical Properties}
\begin{itemize}
\item \textbf{Stability.} Because the learning‑rate $\eta_k$ is inversely proportional to $\sqrt{G_k}$, the algorithm automatically damps large updates when gradients have been large in the past, guaranteeing bounded iterates for $L$‑smooth functions.
\item \textbf{Hyper‑parameter sensitivity.} The global stepsize $\alpha$ appears linearly in the bound; a larger $\alpha$ speeds up early progress but inflates the variance term $L\alpha G_{\max}/\sqrt{T}$. The stability term $\varepsilon$ only affects the first few iterations and does not appear in the asymptotic rate.
\item \textbf{Comparison to SGD.} Standard stochastic gradient descent with a fixed stepsize $\eta$ yields $O(1/\sqrt{T})$ in the non‑convex setting but requires careful tuning of $\eta$. AdaGrad’s adaptive rule replaces $\eta$ by $\alpha/\sqrt{G_k}$, eliminating the need to know $G_{\max}$ a‑priori and providing a data‑dependent stepsize.
\item \textbf{Special case – deterministic gradients.} When $g_k=\nabla f(w_k)$, $G_T=\sum_{k=0}^{T-1}\|\nabla f(w_k)\|^{2}$ and the bound simplifies to 
\[
\frac{1}{T}\sum_{k=0}^{T-1}\|\nabla f(w_k)\|^{2}
\le \frac{2\big(f(w_{0})-f^{\star}\big)}{\alpha\sqrt{T}} .
\]
\end{itemize}

\section*{Preliminaries (Definitions and Lemmas)}
\begin{lemma}[Summation Lemma]\label{lem:summation}
Let $a_1,\dots,a_T\ge0$ and define $S_k:=\sum_{i=1}^{k}a_i$. Then
\[
\sum_{k=1}^{T}\frac{a_k}{\sqrt{S_k}}\le 2\sqrt{\sum_{k=1}^{T}a_k}.
\]
\end{lemma}
\begin{proof}
Consider the telescoping identity
\[
\sqrt{S_k}-\sqrt{S_{k-1}}=\frac{a_k}{\sqrt{S_k}+\sqrt{S_{k-1}}}\ge
\frac{a_k}{2\sqrt{S_k}} .
\]
Summing from $k=1$ to $T$ gives
\[
\sum_{k=1}^{T}\frac{a_k}{\sqrt{S_k}}\le
2\sum_{k=1}^{T}\big(\sqrt{S_k}-\sqrt{S_{k-1}}\big)
=2\sqrt{S_T},
\]
which proves the claim. \qedhere
\end{proof}

\section*{Main Proof (Step‑by‑Step Derivation)}
We prove Theorem~\ref{thm:main}.  
All expectations are taken w.r.t. the randomness of the stochastic gradients up to the current iteration.

\bigskip
\noindent\textbf{[Step 1]} (Smoothness inequality).  
By $L$‑smoothness (Assumption~\ref{ass:smooth}) applied with $y=w_{k+1}=w_k-\eta_k g_k$,
\begin{equation}
f(w_{k+1})\le f(w_k)-\eta_k\langle\nabla f(w_k),g_k\rangle
+\frac{L}{2}\eta_k^{2}\|g_k\|^{2}. \label{eq:smooth}
\end{equation}

\noindent\textbf{[Step 2]} (Unbiasedness).  
Taking conditional expectation $\mathbb{E}[\cdot\mid w_k]$ and using Assumption~\ref{ass:bounded},
\[
\mathbb{E}\big[\langle\nabla f(w_k),g_k\rangle\mid w_k\big]
= \|\nabla f(w_k)\|^{2},\qquad
\mathbb{E}\big[\|g_k\|^{2}\mid w_k\big]\le G_{\max}^{2}.
\]

\noindent\textbf{[Step 3]} (Taking full expectation).  
Apply $\mathbb{E}[\cdot]$ to \eqref{eq:smooth} and use Step~2:
\begin{equation}
\mathbb{E}[f(w_{k+1})]\le
\mathbb{E}[f(w_k)]-\mathbb{E}\big[\eta_k\|\nabla f(w_k)\|^{2}\big]
+\frac{L}{2}\mathbb{E}\big[\eta_k^{2}\|g_k\|^{2}\big]. \label{eq:expect}
\end{equation}

\noindent\textbf{[Step 4]} (Bounding the adaptive stepsizes).  
Recall $\eta_k=\dfrac{\alpha}{\sqrt{G_{k+1}}+\varepsilon}$ with $G_{k+1}=G_k+g_k^{2}$ (element‑wise).  
Since $\varepsilon>0$, for any non‑negative $x$ we have $\frac{1}{x+\varepsilon}\le\frac{1}{x}$, thus
\[
\eta_k\le \frac{\alpha}{\sqrt{G_{k+1}}},\qquad
\eta_k^{2}\le \frac{\alpha^{2}}{G_{k+1}} .
}
\tag{4.1}
\]

\noindent\textbf{[Step 5]} (Insert the bounds).  
Using (4.1) in \eqref{eq:expect} yields
\begin{equation}
\mathbb{E}[f(w_{k+1})]\le
\mathbb{E}[f(w_k)]-
\alpha\;\mathbb{E}\!\left[\frac{\|\nabla f(w_k)\|^{2}}{\sqrt{G_{k+1}}}\right]
+\frac{L\alpha^{2}}{2}\;
\mathbb{E}\!\left[\frac{\|g_k\|^{2}}{G_{k+1}}\right].
\label{eq:bound}
\end{equation}

\noindent\textbf{[Step 6]} (Summation over $k$).  
Sum \eqref{eq:bound} from $k=0$ to $T-1$ and telescope the left‑hand side:
\[
\mathbb{E}[f(w_T)]\le f(w_0)
-\alpha\sum_{k=0}^{T-1}\mathbb{E}\!\left[\frac{\|\nabla f(w_k)\|^{2}}{\sqrt{G_{k+1}}}\right]
+\frac{L\alpha^{2}}{2}\sum_{k=0}^{T-1}
\mathbb{E}\!\left[\frac{\|g_k\|^{2}}{G_{k+1}}\right].
\tag{6.1}
\]

\noindent\textbf{[Step 7]} (Drop the non‑negative term $\mathbb{E}[f(w_T)]\ge f^{\star}$).  
Using Assumption~\ref{ass:lower},
\[
f^{\star}\le\mathbb{E}[f(w_T)] .
\]
Insert this into (6.1) and rearrange:
\begin{equation}
\alpha\sum_{k=0}^{T-1}\mathbb{E}\!\left[\frac{\|\nabla f(w_k)\|^{2}}{\sqrt{G_{k+1}}}\right]
\le f(w_0)-f^{\star}
+\frac{L\alpha^{2}}{2}\sum_{k=0}^{T-1}
\mathbb{E}\!\left[\frac{\|g_k\|^{2}}{G_{k+1}}\right].
\label{eq:prelim}
\end{equation}

\noindent\textbf{[Step 8]} (Bounding the second summation).  
Since $G_{k+1}=G_k+\|g_k\|^{2}\ge \|g_k\|^{2}$ we have
\[
\frac{\|g_k\|^{2}}{G_{k+1}}\le 1 .
\]
Therefore,
\[
\sum_{k=0}^{T-1}\mathbb{E}\!\left[\frac{\|g_k\|^{2}}{G_{k+1}}\right]
\le \sum_{k=0}^{T-1}1 = T .
\tag{8.1}
\]

\noindent\textbf{[Step 9]} (Apply Lemma~\ref{lem:summation}).  
Define $a_k:=\|\nabla f(w_k)\|^{2}$ and $S_k:=\sum_{i=0}^{k} a_i$.
Observe that $G_{k+1}\ge \sum_{i=0}^{k}\|g_i\|^{2}\ge\sum_{i=0}^{k}a_i=S_k$ because $\mathbb{E}\|g_i\|^{2}\ge\|\nabla f(w_i)\|^{2}$ (by Jensen). Hence
\[
\frac{a_k}{\sqrt{G_{k+1}}}\le\frac{a_k}{\sqrt{S_k}} .
\]
Applying Lemma~\ref{lem:summation} to the sequence $\{a_k\}$ gives
\[
\sum_{k=0}^{T-1}\frac{a_k}{\sqrt{G_{k+1}}}
\le
\sum_{k=0}^{T-1}\frac{a_k}{\sqrt{S_k}}
\le 2\sqrt{\sum_{k=0}^{T-1}a_k}
=2\sqrt{\sum_{k=0}^{T-1}\|\nabla f(w_k)\|^{2}} .
\tag{9.1}
\]

\noindent\textbf{[Step 10]} (Combine (9.1) with (7)).  
From \eqref{eq:prelim} and (9.1):
\[
\alpha\;2\sqrt{\sum_{k=0}^{T-1}\mathbb{E}\|\nabla f(w_k)\|^{2}}
\;\le\;
f(w_0)-f^{\star}
+\frac{L\alpha^{2}}{2}\,T .
\]

\noindent\textbf{[Step 11]} (Isolate the sum of gradient norms).  
Divide by $2\alpha$ and square both sides:
\[
\sum_{k=0}^{T-1}\mathbb{E}\|\nabla f(w_k)\|^{2}
\le
\frac{1}{4\alpha^{2}}\Bigl(f(w_0)-f^{\star}
+\frac{L\alpha^{2}}{2}T\Bigr)^{2}.
\tag{11.1}
\]

\noindent\textbf{[Step 12]} (Simplify the right‑hand side).  
Expanding the square and keeping only the dominant terms yields
\[
\sum_{k=0}^{T-1}\mathbb{E}\|\nabla f(w_k)\|^{2}
\le
\frac{(f(w_0)-f^{\star})^{2}}{4\alpha^{2}}
+\frac{L(f(w_0)-f^{\star})}{4}\,T
+\frac{L^{2}\alpha^{2}}{16}\,T^{2}.
\]
Dividing by $T$ gives an average bound:
\[
\frac{1}{T}\sum_{k=0}^{T-1}\mathbb{E}\|\nabla f(w_k)\|^{2}
\le
\frac{(f(w_0)-f^{\star})^{2}}{4\alpha^{2}T}
+\frac{L(f(w_0)-f^{\star})}{4}
+\frac{L^{2}\alpha^{2}}{16}\,T .
\]

\noindent\textbf{[Step 13]} (Use the bounded‑gradient assumption to tighten the variance term).  
From Assumption~\ref{ass:bounded},
$\|g_k\|^{2}\le G_{\max}^{2}$ a.s., hence $G_{k+1}\le (k+1)G_{\max}^{2}$ and consequently
\[
\eta_k=\frac{\alpha}{\sqrt{G_{k+1}}+\varepsilon}
\ge \frac{\alpha}{\sqrt{(k+1)}\,G_{\max}+\varepsilon}
\ge \frac{\alpha}{G_{\max}\sqrt{T}} .
\]
Plugging this lower bound into the left‑hand side of \eqref{eq:prelim} yields
\[
\frac{\alpha}{G_{\max}\sqrt{T}}\sum_{k=0}^{T-1}\mathbb{E}\|\nabla f(w_k)\|^{2}
\le f(w_0)-f^{\star}
+\frac{L\alpha^{2}}{2}T .
\]
Dividing by $T$ finally gives the clean $O(1/\sqrt{T})$ bound:
\begin{equation}
\frac{1}{T}\sum_{k=0}^{T-1}\mathbb{E}\|\nabla f(w_k)\|^{2}
\le
\frac{2\big(f(w_0)-f^{\star}\big)}{\alpha\sqrt{T}}
+\frac{L\alpha G_{\max}}{\sqrt{T}} .
\label{eq:final}
\end{equation}
This is exactly the statement of Theorem~\ref{thm:main}. \qed

\section*{Rate Derivation (With Constants)}
Equation \eqref{eq:final} shows that the average squared gradient norm decays as
\[
\mathcal{O}\!\Bigl(\frac{1}{\sqrt{T}}\Bigr)
\quad\text{with constant}\quad
C = \frac{2\big(f(w_0)-f^{\star}\big)}{\alpha}
+ L\alpha G_{\max}.
\]
Choosing $\alpha=\sqrt{\frac{2\big(f(w_0)-f^{\star}\big)}{L G_{\max}}}$ balances the two terms and yields the optimal constant
\[
C_{\text{opt}} = 2\sqrt{2L G_{\max}\big(f(w_0)-f^{\star}\big)} .
\]

\section*{Special Cases}
\begin{itemize}
\item \textbf{Deterministic gradients} ($g_k=\nabla f(w_k)$).  
Then $G_{\max}= \max_k\|\nabla f(w_k)\|$, and the variance term disappears, leaving
\[
\frac{1}{T}\sum_{k=0}^{T-1}\|\nabla f(w_k)\|^{2}
\le \frac{2\big(f(w_0)-f^{\star}\big)}{\alpha\sqrt{T}} .
\]

\item \textbf{Convex $f$}.  
When $f$ is convex, the same analysis combined with the standard convexity inequality
$f(w_k)-f^{\star}\le\langle\nabla f(w_k), w_k-w^{\star}\rangle$ yields a $O(1/\sqrt{T})$ bound on the suboptimality $f(\bar w_T)-f^{\star}$, matching the classic AdaGrad guarantee.

\item \textbf{Coordinate‑wise AdaGrad}.  
If $d>1$ and the algorithm is applied per coordinate, Lemma~\ref{lem:summation} holds component‑wise, and the final bound becomes a sum over coordinates, still $O(1/\sqrt{T})$.
\end{itemize}

\section*{Discussion}
The proof above follows the classical template for adaptive methods: smoothness gives a one‑step descent inequality, the adaptive stepsize is expressed through the cumulative squared gradients, and Lemma~\ref{lem:summation} converts the resulting telescoping series into a square‑root dependence on the total gradient magnitude. Under the mild bounded‑variance assumption the algorithm enjoys the same $1/\sqrt{T}$ rate as vanilla SGD while automatically adapting to the geometry of the data. The bound is tight for the worst‑case smooth non‑convex setting and improves when gradients are sparse (small $G_{\max}$) or when the loss surface is benign (small $L$).

\end{document}